
\subsubsection{Technical Validation vs. Empirical Findings}

This work validates experimental infrastructure, not empirical properties of code generation models. The key distinction:

\textbf{Infrastructure validated:} Temperature optimization converges, ECE decreases mechanically, accuracy preservation holds mathematically, visualization pipeline functions.

\textbf{Not validated:} Whether Code Llama 7B (or any real code generation model) exhibits high baseline ECE, whether temperature scaling improves real model calibration, whether calibrated confidence correlates with correctness.

\subsubsection{Simulation Design Decisions}

\textbf{Binary logits:} Sufficient to test optimization and metric computation. Real experiments would use full vocabulary-sized logits, but this does not affect the validity of testing whether LBFGS converges or whether ECE can be computed.

\textbf{Overconfidence assumption:} Mock data assumes high baseline ECE (0.527) based on qualitative observations in Kadavath et al. (2022). This assumption is not validated for code generation specifically. Production experiments may reveal different baseline miscalibration levels.

\textbf{No hyperparameter tuning:} The simulation uses standard values from Guo et al. (2017): 15 bins, LBFGS optimizer, initial T=1.5. This avoids overfitting to simulation data and ensures that production experiments use the same protocol.

\subsubsection{Limitations}

\textbf{L1: No real model inference.} This work does not run Code Llama 7B or any other code generation model. All findings are from simulation data designed to test infrastructure, not to represent actual model behavior.

\textbf{L2: Temperature magnitude artifact.} The optimal temperature T* = 2512.71 is not meaningful. Binary logit representation distorts the temperature scale. Production experiments would yield different values.

\textbf{L3: ECE values are simulation artifacts.} Baseline ECE (0.527) and calibrated ECE (0.080) reflect simulation data generation parameters, not properties of real models.

\textbf{L4: No comparison to other calibration methods.} The simulation tests only temperature scaling, not Vector Scaling, Matrix Scaling, or other approaches.

\textbf{L5: Single dataset split.} The simulation uses one data split (200/195). Statistical properties across different splits are not characterized.

\subsubsection{Appropriate Use of These Results}

\textbf{Appropriate:}
\begin{itemize}
\item Citing this work as validation that the experimental protocol is operational
\item Using the validated protocol to design production experiments
\item Referencing the simulation as a test of technical infrastructure
\end{itemize}

\textbf{Inappropriate:}
\begin{itemize}
\item Claiming that code generation models have ECE of 0.53
\item Claiming that temperature scaling reduces ECE by 84.8\% for real models
\item Using simulation results as empirical evidence about model calibration
\end{itemize}

\subsubsection{Future Work}

\textbf{Immediate next step:} Run the validated protocol with real Code Llama 7B on MBPP. This would provide the first empirical measurement of baseline calibration quality for code generation.

\textbf{Methodological extensions:}
\begin{itemize}
\item Compare temperature scaling to Vector Scaling and Matrix Scaling
\item Test calibration across multiple model scales (7B, 13B, 34B)
\item Measure calibration on multiple datasets (HumanEval, APPS, CodeContests)
\end{itemize}

\textbf{Theoretical investigation:}
\begin{itemize}
\item Formalize why (or whether) autoregressive generation amplifies miscalibration
\item Derive theoretical bounds on achievable ECE reduction
\item Investigate relationship between generation length and overconfidence
\end{itemize}

\textbf{Sequential hypothesis validation:}
\begin{itemize}
\item H-M1: Test whether calibrated confidence correlates monotonically with correctness
\item H-M2: Test whether self-critique benefit decreases with confidence
\item H-C1: Test whether confidence-based gating reduces execution attempts
\end{itemize}
